\section{Introduction}\label{sec:introduction}
\subsection{Wildfires, Evacuation and First Nations Peoples}\label{sec:wildfires_and_evacuation}

There are roughly 7,000 wildfires per year in Canada, leading to over 1000 evacuations since 1980. Wildfires are a natural part of the ecosystem but human activity has exacerbated the problem through fires caused directly by humans and by climate change~\citep{cwfis_summary}. They can have devastating impacts on people and the environment through damage to property, ecosystems and loss of life. The study of modelling, predicting and mitigating the impacts of wildfires is an increasingly vital area of research.

While many wildfires occur in unpopulated areas, sometimes they are sufficiently close to humans that evacuation is required. Wildfire evacuation is a complex and challenging process with many competing interests and unique aspects to every fire. For example, the BC Wildfire Service strives to balance saving lives and damage with risk and lives of responders~\citep{bcgov_evacuation_process}. However, evacuation decisions are inevitably subject to the flaws and bias of every institutional decision-making process. 

Canada is home to over 3,000 first nation (FN) reserves~\citep{reservesdb}. Such reserves are close to densely populated cities, but many are remote. Thus, small populations of FN peoples can be particularly vulnerable to risks posed by wildfires. Indeed, FN people consist of just 5\% of the population but, since 1980, 36\% of evacuations have been from FN reserves. Further,~\citet{management_evacuation_review} highlight challenges in decision-making, including limited Indigenous participation and a lack of local control over evacuation processes, which may contribute to inequitable outcomes. Therefore, studying the relationship between wildfires, evacuation and FN peoples is an important area of research to understand the risks and impacts of wildfires on FN peoples, and to inform policies and practices to mitigate these impacts.

This study develops a Bayesian probabilistic model of the evacuation decision by drawing on data from wildfires, evacuations and FN reserves. Section~\ref{sec:introduction} discusses related work and describes the data used. Section~\ref{sec:modelling} introduces the probabilistic models and compares their performance. Section~\ref{sec:results} interprets the results and analyses their implications for wildfire evacuation and FN peoples. Finally, Section~\ref{sec:conclusion} discusses the study's limitations and concludes. All code can be found on Github\footnote{\url{https://github.com/PierreRL/Wildfire_Evacuations}}.

\subsection{Related Work}\label{sec:related_work}

\cite{wildfire_evacs} provide a comprehensive, up-to-date descriptive analysis of wildfire evacuations in Canada, including trends over time, geographic distribution and causes. Their conclusions are broad. With regards to FN peoples, they find that they are disproportionately affected by wildfire evacuations; 14 of the 16 communities in Canada that experienced the most wildfire evacuations were FN. \cite{wildfire_syndromes} analyse over 1000 evacuations by province, local region, time of year and cause to identify three `syndromes' that cover almost 80\% of all evacuations: (1) lightning-caused fires in summer in remote areas, (2) less remote human-caused fires in summer, and (3) interior BC fires in summer. They also analyse evacuations in relation to FN reserves and find that effects and abundance differ across regions. Their work is a useful descriptive analysis of the data, but does not attempt to model the evacuation process. \cite{wildfire_management_evacuation_review} provide a more qualitative analysis of the evacuation process. Their contributions regarding FN peoples focuses on the bias and challenges communities face after evacuation.

A predictive model of evacuation decisions is a useful contribution to the literature. It can be used to identify correlations between wildfire and evacuation characteristics, and to identify potential bias in the evacuation process.

\subsection{Data \& Code}\label{sec:data}

\textbf{Data:} Data on wildfire from the Canadian National Fire Database~\citep{cnfdb} was joined with data on wildfire evacuations from the Canadian Wildfire Evacuation Database~\citep{evacdb}. This join was done based on proximity in space and time. The fire data contains information on the location, size, cause and other characteristics of wildfires. The evacuation data contains information on the location, date and other characteristics of evacuations. Geospatial data on FN reserves from the Canadian Geospatial Data Infrastructure~\citep{reservesdb} was used to compute distance to the nearest FN reserve and whether the fire was in a FN reserve. Categorical variables were simplified and much of the data was dropped. Further details are given in Appendix~\ref{sec:data_appendix} and in the code.

The final dataset contains $N=15,203$ fires, of which 531 had at least one associated evacuation and $269$ were on FN reserves. Predictors are: \texttt{province} (province of fire, $k=12$), \texttt{response\_type} (management response level, $k=4$), \texttt{protection\_zone} (land management category, $k=8$), \texttt{fn\_indicator, binary} (fire in FN reserve), \texttt{log\_dist\_to\_fn\_km} (standard scaled log distance to nearest reserve), \texttt{log\_fire\_size\_ha} (standard scaled log fire size in hectares) and \texttt{region} (categorical variable indicating the grid cell over Canada, defined as squares of approximately $3^\circ \times 3^\circ$ in latitude and longitude, $k=145$ regions with wildfire data). The binary target is \texttt{any\_evacuation} (whether any evacuation occurred).

\textbf{Code:} All models were created in \href{https://mc-stan.org/}{STAN} and compiled and analysed in python using \href{https://github.com/stan-dev/cmdstanpy}{cmdstanpy}. I used 4 chains, 1000 iterations and a warmup of 500 with the default HMC algorithm unless otherwise stated. Analysis of mixing, stability and performance was done using \href{https://github.com/arviz-devs/arviz}{Arviz} and can be found in Section~\ref{sec:modelling}. All code and data can be found on \href{https://github.com/PierreRL/Wildfire_Evacuations}{GitHub} and an example of Stan code used can be found in Appendix~\ref{sec:stan_code}.